% فصل ۳: چالش‌های شبکه اقتضایی خودرویی
% منابع: hozouriOverviewVANETVehicular2023، xuInternetVehiclesBig2018،
% farsimadanReviewSecurityChallenges2025، VanetSecurityPrivacy

\section{چالش‌های شبکه اقتضایی خودرویی}

در مرور چالش‌های شبکه اقتضایی خودرویی، چالش‌ها به سه دستهٔ چالش‌های کاربردها، چالش‌های شبکهٔ داده و چالش‌های مدیریت منابع طبقه‌بندی می‌شوند \cite{hozouriOverviewVANETVehicular2023}. این فصل چالش‌های شبکه‌سازی، مدیریت منابع، مقیاس‌پذیری، استانداردسازی و امنیتی را بررسی می‌کند.

\subsection{چالش‌های شبکه‌سازی}

ویژگی‌های ذاتی شبکه اقتضایی خودرویی، چالش‌هایی در حوزهٔ شبکه‌سازی ایجاد می‌کنند \cite{hozouriOverviewVANETVehicular2023}. کاهش چگالی گره‌ها شبکه را تکه‌تکه می‌کند و تحویل بسته‌ها را مختل می‌سازد و سرعت بالای گره‌ها، توپولوژی را پویاتر و شبکه را چندبخشی می‌کند \cite{hozouriOverviewVANETVehicular2023}. ارتباط میان‌خودرویی \LTRfootnote{\lr{Inter-Vehicular Communication (IVC)}} پایهٔ ارتباطات این شبکه است و مسیریابی داده، امنیت داده و انتشار پیام سه مسئلهٔ بنیادین آن به شمار می‌روند \cite{hozouriOverviewVANETVehicular2023}.

کمبود طیف فرکانسی چالش دیگری است: استانداردهای دسترسی بی‌سیم خودرویی شامل دسترسی بی‌سیم در محیط‌های خودرویی \LTRfootnote{\lr{Wireless Access in Vehicular Environments (WAVE)}} و ارتباطات اختصاصی کوتاه‌برد \LTRfootnote{\lr{Dedicated Short-Range Communications (DSRC)}} هستند؛ این استاندارد ظرفیت محدودی دارد و در شبکه‌های متراکم و بزرگ، مشکلات قابلیت اطمینان و مقیاس‌پذیری در شبکه‌های اقتضایی خودرویی مبتنی بر \lr{DSRC} مشاهده شده است \cite{hozouriOverviewVANETVehicular2023}. هر خودرو به‌طور دوره‌ای از هفت کانال \lr{DSRC} برای تبادل داده استفاده می‌کند و رقابت بر سر این کانال‌ها موجب ازدحام شبکه، افزایش تأخیر بسته و کاهش گذردهی می‌شود \cite{hozouriOverviewVANETVehicular2023}. تأثیرات محیطی نیز بر ارتباط اثر می‌گذارد: ساختمان‌ها، خودروها، درختان و سایر اشیا می‌توانند موجب اختلال شوند و پدیده‌های انتشار چندمسیره، محوشدگی کانال و سایه‌افکنی سیگنال و اثرات آب‌وهوایی را به همراه داشته باشند \cite{hozouriOverviewVANETVehicular2023}.

\subsection{چالش‌های مدیریت منابع}

هم‌زمان با چالش‌های شبکه‌سازی، مدیریت منابع مشترک نیز چالش‌هایی جداگانه ایجاد می‌کند \cite{hozouriOverviewVANETVehicular2023}. در زیرساخت شبکه اقتضایی خودرویی، منابعی چون کانال‌های پهنای باند و واحدهای کنارجاده‌ای \LTRfootnote{\lr{Road-Side Unit (RSU)}} میان خودروها مشترک هستند و استقرار این شبکه به‌دلیل همین اشتراک با چالش‌هایی روبه‌رو می‌شود \cite{hozouriOverviewVANETVehicular2023}. استقرار زیرساخت شبکه اقتضایی خودرویی نیز خود از چالش‌های مدیریت منابع شمرده می‌شود \cite{hozouriOverviewVANETVehicular2023}.

تخصیص منابع رادیویی از دیگر این چالش‌هاست: تکنیک‌های کنترل دسترسی به رسانه \LTRfootnote{\lr{Medium Access Control (MAC)}} برای تضمین دسترسی عادلانهٔ همهٔ خودروها به کانال و کاهش اتلاف بسته ضروری هستند \cite{hozouriOverviewVANETVehicular2023}. پویایی بالا، قطع و وصل گسستهٔ اتصال، نیازهای متنوع کیفیت سرویس \LTRfootnote{\lr{Quality of Service (QoS)}} و دشواری‌های امنیتی از عواملی هستند که این تخصیص را پیچیده می‌کنند \cite{hozouriOverviewVANETVehicular2023}.

\subsection{چالش‌های مقیاس‌پذیری}

در امتداد این چالش‌ها، رشد حجم داده‌ها چالش مقیاس‌پذیری را نیز ایجاد می‌کند \cite{xuInternetVehiclesBig2018}. رشد نمایی داده‌های تولیدشدهٔ خودرویی به‌همراهٔ افزایش تقاضای کاربران درون‌خودرویی، حجم عظیمی از داده در شبکه‌های اقتضایی خودرویی ایجاد کرده است و گسترش هم نوع و هم مقدار داده، چالش‌هایی برای شبکهٔ اقتضایی خودرویی سنتی به همراه دارد \cite{xuInternetVehiclesBig2018}. در سال ۲۰۱۸ پیش‌بینی شده بود که تا سال ۲۰۲۰ یک‌پنجم خودروهای روی جاده به اینترنت متصل شوند و ترافیک خودرویی جهانی به ۳۰۰ هزار اگزابایت برسد \cite{xuInternetVehiclesBig2018}. همچنین پیش‌بینی شده بود که تا سال ۲۰۲۰ بیش از ۲۰۰ حسگر بر روی خودروهای آینده نصب شود و حدود ۴۰۰۰ گیگابایت داده هر روز تولید شود \cite{xuInternetVehiclesBig2018}.

این حجم از داده، به‌ویژه برای پردازش بلادرنگ در رانندگی خودکار، نیازمند ظرفیت ذخیره‌سازی بسیار بزرگ و قابلیت محاسباتی بالا است \cite{xuInternetVehiclesBig2018}. از آنجا که شبکهٔ اینترنت خودروها \LTRfootnote{\lr{Internet of Vehicles (IoV)}} بسیار پویاست، پروتکل‌های شبکه باید نسبت به اندازهٔ شبکه و توپولوژی در حال تغییر مقیاس‌پذیر باشند \cite{xuInternetVehiclesBig2018}. افزونگی داده‌های کلان نیز سربار انتقال، ذخیره‌سازی و پردازش را می‌افزاید و حذف آن ضروری است \cite{xuInternetVehiclesBig2018}.

\subsection{چالش‌های استانداردسازی}

هم‌زمان با این چالش‌های مقیاس، تنوع فناوری‌های دسترسی مسئلهٔ استانداردسازی و تعامل‌پذیری را مطرح می‌کند \cite{farsimadanReviewSecurityChallenges2025}. ارتباطات خودرو با همه‌چیز \LTRfootnote{\lr{Vehicle-to-Everything (V2X)}} بر پایهٔ فناوری‌های گوناگونی چون ارتباطات اختصاصی کوتاه‌برد (\lr{DSRC}/\lr{IEEE 802.11p})، \lr{LTE-V}، \lr{C-V2X} و \lr{5G-NR} عمل می‌کند \cite{farsimadanReviewSecurityChallenges2025}. استاندارد \lr{IEEE 802.11p} که پایهٔ ارتباطات \lr{WAVE} است، تضمین می‌کند که ارتباطات اولویت‌دار بیش از ده‌ها میلی‌ثانیه تأخیر نداشته باشند \cite{farsimadanReviewSecurityChallenges2025}. پهنای باند محدود \lr{DSRC} قادر به پشتیبانی از ارتباطات آیندهٔ \lr{V2X} نیست \cite{farsimadanReviewSecurityChallenges2025} و به همین دلیل استفاده از شبکهٔ ناهمگن \LTRfootnote{\lr{Heterogeneous Network (Het-Net)}} مرکب از \lr{DSRC}، \lr{LTE} و \lr{Wi-Fi} پیشنهاد شده است \cite{farsimadanReviewSecurityChallenges2025}. با افزایش تقاضا برای تأخیر بسیار کم و قابلیت اطمینان بالا، مدل‌های سنتی مدیریت امنیت نیز بدون تحمیل سربار و هزینهٔ عملیاتی قابل‌توجه قادر به پاسخ‌گویی نیستند \cite{farsimadanReviewSecurityChallenges2025}.

تعامل‌پذیری میان ذی‌نفعان چندگانهٔ این حوزه، یعنی اپراتورهای شبکه، مراجع جاده‌ای، خودروسازان (\lr{OEM})، شهرداری‌ها و ارائه‌دهندگان سرویس، چالش دیگری است \cite{farsimadanReviewSecurityChallenges2025}. فناوری‌های دفترکل توزیع‌شده چون بلاکچین می‌توانند تعامل‌پذیری قابل‌اعتماد میان این مشارکت‌کنندگان چندگانه را فراهم کنند، هرچند این امر می‌تواند پیچیدگی زمانی اضافی به همراه داشته باشد \cite{farsimadanReviewSecurityChallenges2025}.

\subsection{چالش‌های امنیتی}

این چالش‌ها همه بر الزامات امنیتی و حریم خصوصی شبکه اقتضایی خودرویی می‌افزایند \cite{farsimadanReviewSecurityChallenges2025, VanetSecurityPrivacy}. برای تضمین امنیت ارتباطات خودرویی، الزاماتی چون دسترس‌پذیری، محرمانگی، تصدیق هویت، یکپارچگی، عدم انکار و پاسخگویی، حریم خصوصی و اعتماد مطرح است \cite{farsimadanReviewSecurityChallenges2025, VanetSecurityPrivacy}. دسترس‌پذیری \LTRfootnote{\lr{Availability}} به‌این معناست که ارتباطات میان اشیاء \lr{V2X} تحلیل و به‌موقع در دسترس گیرندگان مورد نظر قرار گیرند؛ این الزام، استفاده از تکنیک‌های رمزنگاری سبک و کم‌هزینه را ضروری می‌کند و دسترسی حتی در حملات منع سرویس \LTRfootnote{\lr{Denial of Service (DoS)}} نیز باید حفظ شود \cite{farsimadanReviewSecurityChallenges2025}. محرمانگی \LTRfootnote{\lr{Confidentiality}} مستلزم رمزنگاری داده‌هاست؛ هرچند اطلاعات تبادل‌شده در \lr{V2X} به‌طور کلی عمومی است و محرمانگی ممکن است جزو اولویت‌های بالا نباشد، جز در مورد اطلاعات مربوط به حریم خصوصی سرنشینان خودرو \cite{farsimadanReviewSecurityChallenges2025}. تصدیق هویت \LTRfootnote{\lr{Authentication}} امکان تمایز میان اشیاء معتبر و عناصر موذی را فراهم می‌کند و دو زیرنوع دارد: تصدیق پیام، یعنی اصالت و عدم دستکاری بسته، و تصدیق کاربر یا منبع، یعنی اعتبار موجودیت \cite{farsimadanReviewSecurityChallenges2025}. یکپارچگی \LTRfootnote{\lr{Integrity}} مستلزم آن است که داده‌ها به‌موقع بررسی شوند تا هرگونه دستکاری، تغییر یا حذف در حین انتقال شناسایی گردد \cite{farsimadanReviewSecurityChallenges2025}.

عدم انکار \LTRfootnote{\lr{Non-repudiation}} و پاسخگویی \LTRfootnote{\lr{Accountability}} به اثبات ارتباط و تراکنش‌های داده برای جلوگیری از انکار مشارکت طرفین مربوط است و ابزارهای آن شامل امضای دیجیتال، \lr{PKI} \LTRfootnote{\lr{Public Key Infrastructure (PKI)}}، مهر زمانی و ثبت وقایع امن است؛ این الزام کاربرد حقوقی پس از تصادف دارد \cite{VanetSecurityPrivacy}. پاسخگویی در چارچوب حریم خصوصی نیز اهمیت دارد: حریم خصوصی باید اطلاعات شخصی و مکانی را در برابر ردیابی و نمایه‌سازی محافظت کند، در حالی که هم‌زمان امکان پاسخ‌گویی در اختلافات یا تحقیقات را حفظ می‌کند \cite{VanetSecurityPrivacy}.

حریم خصوصی \LTRfootnote{\lr{Privacy}} در این شبکه گره‌ای است که گمنامی زیرمجموعه‌ای از آن شمرده می‌شود؛ واحدهای کنارجاده‌ای می‌توانند در فرایند تصدیق، موقعیت خودرو را ردیابی کنند، ردیابی‌ناپذیری به این معناست که فعالیت‌های خودرو قابل پیگیری نباشد و پیوندناپذیری به این معنا که طرف غیرمجاز نتواند هویت خودرو را به راننده یا مالک آن پیوند زند \cite{farsimadanReviewSecurityChallenges2025}. اعتماد \LTRfootnote{\lr{Trust}} نیز به‌عنوان معیاری تکمیلی تعریف می‌شود: اطمینان یک موجودیت نسبت به موجودیت دیگری که جزئی از شبکهٔ خودرویی است، و محاسبهٔ اعتماد معیار امنیتی اضافی محسوب می‌شود \cite{farsimadanReviewSecurityChallenges2025}.
